% Unit 1: Foundations and Prerequisites
% Master Unit Reference Guide (Cheat Sheet, Problems, and Solutions)
% Prepared by Staples Education
% Content informed by the local curriculum modules (1.1 to 1.6) and the
% QUIZ_DATA micro-practice and unit-mastery items for Unit 1.

\documentclass[11pt]{article}

\usepackage[margin=1in]{geometry}
\usepackage{amsmath}
\usepackage{amssymb}
\usepackage{xcolor}
\usepackage[most]{tcolorbox}
\usepackage{enumitem}
\usepackage{hyperref}
\hypersetup{colorlinks=true, linkcolor=headerblue, urlcolor=headerblue}

\tcbuselibrary{skins}

% ---------------------------------------------------------------------------
% Unified style-guide palette (see LATEX_STYLE_GUIDE.md)
% ---------------------------------------------------------------------------
\definecolor{headerblue}{HTML}{1C569E}   % overview / structural framing
\definecolor{accentgreen}{HTML}{2E7D32}  % algorithms / methods
\definecolor{accentorange}{HTML}{E27A1E} % formulas / concept takeaways
\definecolor{workpurple}{HTML}{A020F0}   % worked-example tracking (vibrant)
\definecolor{warnred}{HTML}{C0392B}      % crimson pitfall / warning

% ---------------------------------------------------------------------------
% Subsection typography: headerblue numbered sub-module titles
% ---------------------------------------------------------------------------
\usepackage{titlesec}
\titleformat{\subsection}
  {\normalfont\large\bfseries\color{headerblue}}{\thesubsection}{1em}{}

% ---------------------------------------------------------------------------
% Six core custom boxes (unbreakable, title={#1} protected)
% ---------------------------------------------------------------------------

% 1. Blue overview capsule for the topics summary.
\newtcolorbox{topicsbox}{
    enhanced,
    colback=headerblue!6!white, colframe=headerblue,
    coltitle=white, fonttitle=\bfseries,
    title=Unit 1 at a Glance: Core Sub-Topics,
    boxrule=0.6pt, arc=2mm,
    left=3mm, right=3mm, top=2mm, bottom=2mm
}

% 2. Orange formula container for standard forms and master formulas.
\newtcolorbox{formulabox}[1][Master Formula]{
    enhanced,
    colback=accentorange!8!white, colframe=accentorange,
    coltitle=white, fonttitle=\bfseries,
    title={#1},
    boxrule=0.6pt, arc=1mm,
    left=3mm, right=3mm, top=2mm, bottom=2mm
}

% 3. Green algorithm block for explicitly laid-out procedures.
\newtcolorbox{methodbox}[1][Method]{
    enhanced,
    colback=accentgreen!6!white, colframe=accentgreen,
    coltitle=white, fonttitle=\bfseries,
    title={#1},
    boxrule=0.6pt, arc=1mm,
    left=3mm, right=3mm, top=2mm, bottom=2mm
}

% 4. Purple west-bordered worked example for step-by-step solutions.
\newtcolorbox{examplebox}[1][Worked Example]{
    enhanced,
    colback=workpurple!4!white, colframe=workpurple,
    coltitle=white, fonttitle=\bfseries,
    title={#1},
    boxrule=0.4pt, sharp corners,
    borderline west={3pt}{0pt}{workpurple},
    left=5mm, right=3mm, top=2mm, bottom=2mm
}

% 5. Orange thin-bordered final concept takeaway.
\newtcolorbox{conceptbox}[1][Key Takeaway]{
    enhanced,
    colback=accentorange!5!white, colframe=accentorange,
    coltitle=white, fonttitle=\bfseries,
    title={#1},
    boxrule=0.3pt, arc=1mm,
    left=3mm, right=3mm, top=2mm, bottom=2mm
}

% 6. Crimson west-bordered warning box for common pitfalls.
\newtcolorbox{warningbox}[1][Common Pitfall]{
    enhanced,
    colback=red!3!white, colframe=warnred,
    coltitle=white, fonttitle=\bfseries,
    title={#1},
    boxrule=0.4pt, sharp corners,
    borderline west={3pt}{0pt}{warnred},
    left=5mm, right=3mm, top=2mm, bottom=2mm
}

\setlist[enumerate]{itemsep=4pt, topsep=4pt}

% Number and index sections plus subsections (e.g., 1.1) two levels deep.
\setcounter{secnumdepth}{2}
\setcounter{tocdepth}{2}

\begin{document}

% ---------------------------------------------------------------------------
% Title block
% ---------------------------------------------------------------------------
\begin{center}
    {\LARGE \bfseries Unit 1: Foundations and Prerequisites}\\[6pt]
    {\large Prepared by Staples Education}\\[4pt]
    {\itshape Master Unit Reference Guide: Cheat Sheet, Problems, and Solutions}
\end{center}

\vspace{0.4em}

% ---------------------------------------------------------------------------
% Table of contents (master guide)
% ---------------------------------------------------------------------------
\tableofcontents

\vspace{0.4em}

% ---------------------------------------------------------------------------
% Topics overview box
% ---------------------------------------------------------------------------
\begin{topicsbox}
This unit refreshes the calculus toolkit and the vocabulary the rest of the
course assumes. It moves from derivatives and integrals, through the number $e$
and Euler's formula, into the precise language of order, degree, linearity, and
what it means to be a solution.
\begin{itemize}[leftmargin=1.4em, itemsep=3pt]
    \item \textbf{Calculus Review} --- the power, product, and chain rules, and
          integrating the exponential.
    \item \textbf{The Number $e$} --- the function that is its own derivative,
          and where it comes from.
    \item \textbf{Euler's Formula and Identity} --- $e^{i\theta} = \cos\theta +
          i\sin\theta$ and the geometry of rotation.
    \item \textbf{Order, Degree, and Linearity} --- the classification words
          that label every equation.
    \item \textbf{What It Means to Be a Solution} --- verification by
          substitution, general versus particular.
    \item \textbf{Eliminating Arbitrary Constants} --- building a differential
          equation from a family of curves.
\end{itemize}
\end{topicsbox}

\vspace{0.6em}
\hrule
\vspace{0.6em}

% ===========================================================================
% PART I: CONDENSED CHEAT SHEET
% ===========================================================================
\phantomsection
\addcontentsline{toc}{section}{Part I --- Condensed Cheat Sheet}
\section*{Part I --- Condensed Cheat Sheet}
\setcounter{section}{1}   % subsections auto-number as the unit modules 1.1, 1.2, ...

\subsection{Calculus Review}

Every solving technique ahead rests on fluent differentiation and integration.
The three rules you reach for most are the power rule, the product rule, and the
chain rule --- and reversing them is exactly what integration asks of you.

\begin{formulabox}[Core Derivative and Integral Rules]
\[
    \frac{d}{dx}x^{n} = nx^{n-1}, \qquad
    \frac{d}{dx}\big[f g\big] = f'g + fg', \qquad
    \frac{d}{dx}f\big(g(x)\big) = f'\big(g(x)\big)\,g'(x).
\]
Integrating the exponential reverses the chain-rule factor:
\[
    \int e^{kx}\,dx = \frac{1}{k}e^{kx} + C.
\]
\end{formulabox}

\begin{warningbox}[Common Pitfall: The Missing Inner Derivative and the Lost $+C$]
The chain rule multiplies by the derivative of the inside --- $\frac{d}{dx}e^{3x}
= 3e^{3x}$, not $e^{3x}$. And an indefinite integral always carries the
arbitrary constant $+C$; dropping it discards the entire family of solutions an
initial condition is meant to select from.
\end{warningbox}

\begin{examplebox}[Worked Example 1.1: Differentiate two composites]
\textbf{(a)} Chain rule on $e^{3x}$:
\[ \frac{d}{dx}e^{3x} = 3e^{3x}. \]
\textbf{(b)} Power and chain rule on $(2x+1)^{3}$:
\[ \frac{d}{dx}(2x+1)^{3} = 3(2x+1)^{2}\cdot 2 = 6(2x+1)^{2}. \]
\end{examplebox}

\subsection{The Number $e$}

The natural exponential $e^{x}$ is singled out by one property: it is its own
derivative. That self-matching behavior is exactly what a rate proportional to
the current amount demands, which is why $e$ governs every growth and decay
model in the course.

\begin{formulabox}[The Natural Exponential]
\[
    \frac{d}{dx}e^{x} = e^{x}, \qquad
    \frac{d}{dt}e^{kt} = k\,e^{kt}.
\]
The constant itself arises as a limit and as a series:
\[
    e = \lim_{n\to\infty}\left(1 + \frac{1}{n}\right)^{n}
      = \sum_{n=0}^{\infty}\frac{1}{n!} \approx 2.718.
\]
\end{formulabox}

\begin{examplebox}[Worked Example 1.2: Solve $\frac{dy}{dt} = y$]
The equation asks for a function equal to its own derivative. The natural
exponential matches, scaled by a free constant:
\[ y = Ce^{t}, \qquad \frac{dy}{dt} = Ce^{t} = y. \quad\checkmark \]
The constant $C$ supplies the whole family of solutions.
\end{examplebox}

\subsection{Euler's Formula and Identity}

Euler's formula ties the exponential to rotation: feeding an imaginary number
into $e^{x}$ traces the unit circle. This is the bridge that lets a single
complex exponential stand in for an oscillating sine--cosine pair.

\begin{formulabox}[Euler's Formula, Identity, and Polar Form]
\[
    e^{i\theta} = \cos\theta + i\sin\theta, \qquad
    e^{i\pi} + 1 = 0.
\]
A complex number in polar form is $z = re^{i\theta}$, where $r = |z|$ is the
distance from the origin. Every point on the unit circle has $\big|e^{i\theta}\big| = 1$,
and products multiply moduli while adding angles.
\end{formulabox}

\begin{examplebox}[Worked Example 1.3: Evaluate by rotation]
\textbf{(a)} A quarter-turn:
\[ e^{i\pi/2} = \cos\tfrac{\pi}{2} + i\sin\tfrac{\pi}{2} = i. \]
\textbf{(b)} Multiply in polar form, multiplying moduli and adding angles:
\[ \big(2e^{i\pi/6}\big)\big(3e^{i\pi/3}\big) = 6\,e^{i(\pi/6 + \pi/3)} = 6\,e^{i\pi/2}. \]
\end{examplebox}

\subsection{Basic Definitions: Order, Degree, and Linearity}

Before any technique applies, you classify the equation. The \emph{order} is the
highest derivative present; the \emph{degree} is the power on that highest
derivative once the equation is polynomial in its derivatives; and
\emph{linearity} asks whether the unknown and its derivatives appear only to the
first power.

\begin{formulabox}[Classification Vocabulary]
\begin{itemize}[leftmargin=1.4em, itemsep=2pt]
    \item \textbf{Order} --- the highest derivative that appears.
    \item \textbf{Degree} --- the power on the highest-order derivative.
    \item \textbf{Linear} --- $y$ and its derivatives appear to the first power
          only, never squared and never multiplied together.
    \item \textbf{ODE vs PDE} --- one independent variable versus several.
    \item An $n$-th order equation needs $n$ initial conditions to pin a unique
          solution.
\end{itemize}
\end{formulabox}

\begin{warningbox}[Common Pitfall: Order Is Not Degree]
Order and degree are different questions. In $\left(\dfrac{dy}{dx}\right)^{3} + y
= x$ the highest derivative is the \emph{first}, so the order is one, while the
cube makes the degree three. And $\dfrac{dy}{dx} + y^{2} = 0$ is
\emph{nonlinear} --- the squared $y$, not the derivative, breaks the rule.
\end{warningbox}

\begin{examplebox}[Worked Example 1.4: Classify three equations]
\textbf{(a)} $\dfrac{d^{2}y}{dx^{2}} + 3\dfrac{dy}{dx} = 0$ --- order $2$,
degree $1$, linear.\\
\textbf{(b)} $\left(\dfrac{dy}{dx}\right)^{3} + y = x$ --- order $1$, degree $3$,
nonlinear.\\
\textbf{(c)} $\dfrac{d^{3}y}{dx^{3}} + \left(\dfrac{dy}{dx}\right)^{2} = 0$ ---
order $3$, degree $1$ (the power sits on the first derivative, not the highest).
\end{examplebox}

\subsection{What It Means to Be a Solution}

A solution is a function that makes the equation an identity --- true for every
$x$, not at one lucky point. You confirm a candidate by substituting it and its
derivatives and checking both sides agree. The \emph{general} solution keeps its
arbitrary constants free; a \emph{particular} solution fixes them.

\begin{methodbox}[Algorithm: Verifying a Candidate Solution]
\begin{enumerate}[leftmargin=1.6em, itemsep=2pt]
    \item Compute every derivative of the candidate the equation requires.
    \item Substitute the candidate and those derivatives into the equation.
    \item Simplify the left and right sides independently.
    \item Confirm the two sides are equal as an identity in $x$ --- for all
          inputs, not a single point.
\end{enumerate}
\end{methodbox}

\begin{warningbox}[Common Pitfall: ``It Works at One Point'']
Checking a candidate at $x = 0$ alone proves nothing --- almost any function
agrees somewhere. A genuine solution must satisfy the equation across its
\emph{whole} domain.
\end{warningbox}

\begin{examplebox}[Worked Example 1.5: Verify two candidates]
\textbf{(a)} Does $y = e^{2x}$ solve $y'' - 4y = 0$? Here $y'' = 4e^{2x}$ and
$4y = 4e^{2x}$, so $y'' - 4y = 0$. \checkmark\\
\textbf{(b)} Does $y = \cos x$ solve $y'' + y = 0$? Here $y'' = -\cos x$, so
$y'' + y = -\cos x + \cos x = 0$. \checkmark
\end{examplebox}

\subsection{Forming Differential Equations by Eliminating Constants}

The reverse problem is just as instructive: given a family of curves with $n$
arbitrary constants, you build the differential equation it satisfies by
differentiating $n$ times and eliminating the constants. The resulting order
equals the number of independent constants.

\begin{methodbox}[Algorithm: Eliminating Arbitrary Constants]
\begin{enumerate}[leftmargin=1.6em, itemsep=2pt]
    \item Count the independent arbitrary constants; call it $n$.
    \item Differentiate the relation $n$ times to produce $n$ new equations.
    \item Use those equations to solve for and eliminate every constant.
    \item The result is an $n$-th order differential equation, free of constants.
\end{enumerate}
\end{methodbox}

\begin{examplebox}[Worked Example 1.6: Eliminate the constants]
\textbf{(a)} From $y = Cx$: differentiate to get $y' = C$, and since $C = y/x$,
\[ \frac{dy}{dx} = \frac{y}{x}. \]
\textbf{(b)} From $y = A\cos x + B\sin x$ (two constants): differentiate twice,
$y'' = -A\cos x - B\sin x = -y$, so
\[ \frac{d^{2}y}{dx^{2}} + y = 0. \]
\end{examplebox}

% ===========================================================================
% PART II: REVIEW GUIDE (PRACTICE QUESTIONS)
% ===========================================================================
\clearpage
\phantomsection
\addcontentsline{toc}{section}{Part II --- Review Guide: Practice Problems}
\section*{Part II --- Review Guide: Practice Problems}

Work the following ten problems in order. They are graded from raw calculus
review, through the number $e$ and Euler's formula, into classifying equations,
verifying solutions, and finally building equations by eliminating constants.
Complete solutions follow in Part III.

\begin{enumerate}[leftmargin=2em]
    \item Differentiate $y = e^{3x}$ using the chain rule.

    \item Evaluate the indefinite integral $\displaystyle\int e^{-2x}\,dx$.

    \item Differentiate \textbf{(a)} $(2x+1)^{3}$ and \textbf{(b)} $\sin(x^{2})$.

    \item Compute $\dfrac{d}{dt}e^{kt}$, then state the general solution of
          $\dfrac{dy}{dt} = y$.

    \item Use Euler's formula $e^{i\theta} = \cos\theta + i\sin\theta$ to
          evaluate \textbf{(a)} $e^{i\pi}$ and \textbf{(b)} $e^{i\pi/2}$, and
          write down Euler's identity.

    \item For each equation, state the order and the degree, and say whether it
          is linear:
          \textbf{(a)} $\dfrac{d^{2}y}{dx^{2}} + 3\dfrac{dy}{dx} = 0$, \quad
          \textbf{(b)} $\left(\dfrac{dy}{dx}\right)^{3} + y = x$, \quad
          \textbf{(c)} $\dfrac{dy}{dx} + y^{2} = 0$.

    \item Verify that $y = e^{2x}$ is a solution of $\dfrac{d^{2}y}{dx^{2}} - 4y = 0$.

    \item Decide whether $y = x^{2}$ satisfies $x\dfrac{dy}{dx} = 2y$, and
          whether $y = \cos x$ satisfies $\dfrac{d^{2}y}{dx^{2}} + y = 0$.

    \item Form a differential equation by eliminating $C$ from \textbf{(a)}
          $y = Cx^{2}$ and \textbf{(b)} $y = Ce^{x}$.

    \item \textbf{(Capstone)} Eliminate both constants from
          $y = A\cos x + B\sin x$, and separately eliminate $r$ from the circle
          family $x^{2} + y^{2} = r^{2}$.
\end{enumerate}

% ===========================================================================
% PART III: SOLUTIONS
% ===========================================================================
\clearpage
\phantomsection
\addcontentsline{toc}{section}{Part III --- Complete Solutions}
\section*{Part III --- Complete Solutions}

\noindent\textbf{Solution 1.}\quad The exponential reproduces itself, and the
chain rule multiplies by the derivative of the inside, $3$.
\[ \textcolor{workpurple}{\frac{dy}{dx} = e^{3x}\cdot 3 = 3e^{3x}}. \]

\medskip
\noindent\textbf{Solution 2.}\quad Integrating an exponential divides by the
constant in the exponent and restores the arbitrary constant.
\[ \textcolor{workpurple}{\int e^{-2x}\,dx = \frac{1}{-2}e^{-2x} + C = -\tfrac{1}{2}e^{-2x} + C}. \]

\medskip
\noindent\textbf{Solution 3.}\quad Apply the power-and-chain rule, then the
chain rule on the sine.
\[ \textcolor{workpurple}{\frac{d}{dx}(2x+1)^{3} = 3(2x+1)^{2}\cdot 2 = 6(2x+1)^{2}} \]
\[ \textcolor{workpurple}{\frac{d}{dx}\sin(x^{2}) = \cos(x^{2})\cdot 2x = 2x\cos(x^{2})}. \]

\medskip
\noindent\textbf{Solution 4.}\quad The exponential returns itself times the
constant from the exponent.
\[ \textcolor{workpurple}{\frac{d}{dt}e^{kt} = k\,e^{kt}}. \]
The equation $\frac{dy}{dt} = y$ asks for a function equal to its own derivative,
so
\[ \textcolor{workpurple}{y = Ce^{t}}. \]

\medskip
\noindent\textbf{Solution 5.}\quad Substitute the angles into Euler's formula.
\[ \textcolor{workpurple}{e^{i\pi} = \cos\pi + i\sin\pi = -1 + 0 = -1} \]
\[ \textcolor{workpurple}{e^{i\pi/2} = \cos\tfrac{\pi}{2} + i\sin\tfrac{\pi}{2} = 0 + i = i}. \]
Rearranging the first gives Euler's identity,
$\textcolor{workpurple}{e^{i\pi} + 1 = 0}$.

\medskip
\noindent\textbf{Solution 6.}\quad The order is the highest derivative; the
degree is its power; linearity forbids powers above one on $y$ and its
derivatives.
\[ \textcolor{workpurple}{\text{(a) order } 2,\; \text{degree } 1,\; \text{linear}} \]
\[ \textcolor{workpurple}{\text{(b) order } 1,\; \text{degree } 3,\; \text{nonlinear}} \]
\[ \textcolor{workpurple}{\text{(c) order } 1,\; \text{degree } 1,\; \text{nonlinear (the } y^{2} \text{ term)}}. \]

\medskip
\noindent\textbf{Solution 7.}\quad Differentiate twice and substitute.
\[ \textcolor{workpurple}{\frac{dy}{dx} = 2e^{2x}}, \qquad
   \textcolor{workpurple}{\frac{d^{2}y}{dx^{2}} = 4e^{2x}} \]
\[ \textcolor{workpurple}{\frac{d^{2}y}{dx^{2}} - 4y = 4e^{2x} - 4e^{2x} = 0} \qquad\checkmark \]
so $y = e^{2x}$ is a solution.

\medskip
\noindent\textbf{Solution 8.}\quad Test each candidate by substitution.
\[ \textcolor{workpurple}{y = x^{2}: \quad y' = 2x, \;\; x\,y' = x\cdot 2x = 2x^{2} = 2y} \qquad\checkmark \]
\[ \textcolor{workpurple}{y = \cos x: \quad y'' = -\cos x, \;\; y'' + y = -\cos x + \cos x = 0} \qquad\checkmark \]
Both candidates check out.

\medskip
\noindent\textbf{Solution 9.}\quad Differentiate, then replace the constant from
the original relation.
\[ \textcolor{workpurple}{y = Cx^{2}: \quad y' = 2Cx, \;\; C = \frac{y}{x^{2}} \;\Longrightarrow\; y' = 2\cdot\frac{y}{x^{2}}\cdot x = \frac{2y}{x}} \]
\[ \textcolor{workpurple}{\text{so } x\frac{dy}{dx} = 2y}. \]
\[ \textcolor{workpurple}{y = Ce^{x}: \quad y' = Ce^{x} = y \;\Longrightarrow\; \frac{dy}{dx} = y}. \]

\medskip
\noindent\textbf{Solution 10.}\quad Each family needs differentiation to clear
its constants. For $y = A\cos x + B\sin x$ (two constants), differentiate twice:
\[ \textcolor{workpurple}{y'' = -A\cos x - B\sin x = -y \;\Longrightarrow\; \frac{d^{2}y}{dx^{2}} + y = 0}. \]
For the circle $x^{2} + y^{2} = r^{2}$ (one constant), differentiate implicitly:
\[ \textcolor{workpurple}{2x + 2y\frac{dy}{dx} = 0 \;\Longrightarrow\; \frac{dy}{dx} = -\frac{x}{y}}. \]

\vspace{0.8em}
\begin{conceptbox}[Unit 1 Key Takeaway]
This unit is the vocabulary and toolkit the whole course speaks in, exercised in
both directions. Going forward, you \textcolor{accentorange}{differentiate and
integrate} fluently and \textcolor{accentorange}{verify a solution} by
substituting it into the equation; going backward, you
\textcolor{accentorange}{eliminate arbitrary constants} to build the equation a
family satisfies. The \textcolor{accentorange}{natural exponential $e$} is the
function equal to its own derivative, and the
\textcolor{accentorange}{order, degree, and linearity} of an equation are the
first labels you read off before any method is chosen.
\end{conceptbox}

\end{document}
